\documentclass[11pt,a4paper]{article}
\usepackage{fontspec}
\usepackage{kotex}
\setmainfont{Times New Roman}
\setmainhangulfont{Batang}
\setsanshangulfont{Malgun Gothic}
\usepackage[margin=25mm,headheight=15pt]{geometry}
\usepackage{amsmath,amssymb,booktabs,array,longtable,graphicx}
\usepackage{setspace,placeins,fancyhdr,xurl,hyperref}
\hypersetup{hidelinks,pdftitle={이주배경 인구가 집중된 지역에서 아동보호 행정 대응은 더 느린가},pdfauthor={}}
\setstretch{1.36}
\setlength{\parindent}{1em}
\setlength{\parskip}{3pt}
\setlength{\tabcolsep}{3pt}
\renewcommand{\arraystretch}{1.22}
\renewcommand{\thesection}{\Roman{section}}
\renewcommand{\thesubsection}{\arabic{subsection}}
\renewcommand{\thesubsubsection}{\thesubsection.\arabic{subsubsection}}
\renewcommand{\tablename}{표}
\renewcommand{\refname}{참고문헌}
\renewcommand{\abstractname}{국문초록}
\pagestyle{fancy}\fancyhf{}\fancyhead[L]{\small 스웨덴 아동보호 행정의 두 사례}\fancyhead[R]{\small\thepage}\renewcommand{\headrulewidth}{0pt}
\begin{document}
\begin{center}
{\LARGE\bfseries 이주배경 인구가 집중된 지역에서\\[5pt] 아동보호 행정 대응은 더 느린가}\par\vspace{10pt}
{\large 스웨덴 지방정부의 사회복지 조사기간과 학업보장 절차에 관한 두 사례 분석}\par\vspace{8pt}
{\small 통합 연구원고\quad 2026년 11월 18일}
\end{center}
\begin{abstract}
이주배경 인구가 집중된 지역에서 아동보호 행정 대응이 더 느린가를 두 가지 행정지표로 살펴본다. CASE 1은 스웨덴 51개 대도시권 지방정부의 일반 아동·청소년 사회복지 조사기간을, CASE 2는 전국 지방정부의 학업보장 서면 절차 충족 응답을 분석한다. 두 자료의 표본과 추정량을 구분하였다. CASE 1의 2024년 횡단면 관계는 작았으나, 지방정부·연도 고정효과 모형에서 스웨덴 출생 아동 중 양부모 해외 출생 비중의 계수는 1\%포인트당 6.70일이었다. 복수 영향 지역을 제외하거나 차분 구간을 바꾸면 추론이 달라져 일관된 기간 증가로 일반화하기 어렵다. 공식 원자료로 재구축한 CASE 2의 191개 장기패널에서는 고비중 집단의 충족률이 2023·2025·2026년 모두 비교집단보다 높았고, 격차는 5.48, 20.78, 3.66\%포인트였다. 독립 DML-DID의 2026--2025년 주 추정은 $-18.43$\%포인트였으며 95\% 구간은 $[-40.09,3.22]$였다. 학습기 난수를 고정하고 주별 폴드만 바꾼 10개 설정에서는 모두 음수이고 8개가 5\% 수준에서 유의했지만, 이를 하나의 통합 검정으로 해석하지 않았다. 2026--2023년 변화 차이는 이 점검의 어느 설정에서도 유의하지 않았다. 263개 연결표본에서도 단기 격차는 축소됐으나 고비중 집단의 충족 건수는 17개에서 31개로 늘었다가 27개가 되었고, 비교집단은 39개, 56개, 69개로 증가했다. 세 시점의 충족률은 고비중 집단에서 모두 더 높았다. 따라서 2025--2026년의 음의 변화 차이는 고비중 집단의 충족률이 비교집단보다 낮아졌다는 뜻이 아니다. 2025년에 커졌던 양의 격차가 줄었다는 뜻이다. 이 수준과 변화의 구분이 뒤의 DML 결과를 해석하는 출발점이다.

\end{abstract}
\noindent\textbf{주제어:} 이주배경, 아동보호 행정, 사회복지 조사기간, VIND28, DML-DID, 응답 전환\par
\noindent\textbf{JEL:} C23, I21, I38, J15

\section{서론}
\subsection{연구 배경과 문제 제기}
이주배경 인구가 많은 지역의 아동보호 행정은 더 느린가. 언어지원과 기관 간 조정이 필요한 사건의 구성은 지역마다 다르고, 담당 인력과 축적된 경험도 같지 않다. 지원 수요가 늘어 조사가 길어질 수도 있지만, 다양한 주민을 지원해 온 지역의 전문성과 절차가 그 부담을 줄일 수도 있다. 인구구성만으로 행정 대응의 방향을 정할 수 없는 이유다.

이 질문을 확인하려면 무엇이 느린지를 먼저 정해야 한다. 일반 사회복지 조사를 끝내기까지의 시간, 필요한 서면 절차의 충족, 아동에게 실제 지원이 시작된 시점은 서로 다른 결과다. 조사가 진행되는 동안 긴급지원이 먼저 제공될 수 있고, 서면 절차가 있어도 개별 아동의 등교에는 시간이 걸릴 수 있다. 이 글에서 처리 속도를 직접 살펴보는 지표는 CASE 1의 평균 조사기간이다. CASE 2의 VIND28은 보호시설 동반 아동의 학업보장 절차에 관한 응답으로, 실제 대기시간을 대신하지 않는다.

두 사례는 서로 다른 방식으로 단순한 지역 비교를 다시 보게 한다. 조사기간은 2024년의 지역 간 비교와 같은 지역의 시간 변화에서 결과가 달랐다. 절차 응답은 고비중 집단이 비교집단보다 높은 충족률을 보이면서도, 2025년에 벌어졌던 격차가 2026년에 줄었다. 이 축소가 고비중 지역의 후퇴인지 비교지역의 개선인지, 어느 표본과 추정 설정에서도 같은 설명이 가능한지를 구분할 필요가 있다. 이 글은 느린 대응을 전제하지 않고 그 명제를 지지하는 증거의 범위를 검토한다.

\subsection{연구 질문 및 기존 문헌과의 차별성}
중심 질문은 이주배경 인구가 집중된 지역에서 아동보호 행정 대응이 더 느린가이다. RQ1은 인구비중이 높은 지방정부에서 일반 아동·청소년 조사의 평균 기간이 더 긴지, RQ2는 사전 인구구성으로 구분한 집단 사이에서 학업보장 절차의 충족률과 변화 경로가 다른지, RQ3는 그 관계가 관측기간·표본·모형을 바꾸어도 유지되는지를 묻는다. CASE 2의 시군별 응답 전환은 동일한 아동이나 실제 절차 운영을 추적한 기록과 구분한다.

이 글의 기여는 여러 방법의 유의한 계수를 모으는 데 있지 않다. 첫째, 처리기간과 서면 절차를 서로 다른 행정 결과로 측정한다. 둘째, 절차 충족률의 수준과 변화, 집단 격차와 개별 지역의 응답 전환을 연결해 격차 축소의 구성을 밝힌다. 셋째, 영향 지역 제외와 표본 확장, DML의 폴드 배정 점검을 통해 추론이 달라지는 지점을 제시한다. 두 사례의 원고와 공개 비교 지표에 관한 연구를 바탕으로\cite{case1,case2,trygged,denvall}, 집계 지표가 말해 주는 행정 상태와 이용 아동의 경험을 구분한다. 두 결과를 합친 회귀나 통합 효과는 추정하지 않는다.

\subsection{연구의 실증적 접근과 구성}
CASE 1은 스톡홀름·예테보리·말뫼 대도시권 51개 지방정부의 Kolada와 SCB 자료를 사용한다. 2024년 단면과 2022--2025년 패널을 구분하며 아동 배경변수는 실제 값이 있는 2022--2024년만 사용한다. CASE 2는 VIND28의 191개 장기 균형패널과 263개 2개년 연결패널을 구분한다. 일반 DID를 집단 변화 차이의 기준선으로 두고 DR-DID·DML-DID·SDID와 측정·표본 민감도를 함께 읽는다.

이 통합은 제공된 두 PDF의 결과와 기존 CASE 1 계산 파일을 대조하여 작성하였다. CASE 2는 Socialstyrelsen의 공식 2026년 Excel과 SCB의 2015년 원자료를 새로 내려받아 패널을 재구축하고 DML-DID를 독립 재추정하였다. 원래 실행 코드는 발견하지 못했으므로 기존 구현을 그대로 복원한 정확 재현과는 구분한다. 기존 DR-DID·SDID와 추가 원고 보고치는 새 추정치와 별도로 표시한다. 이후 II절은 공통 배경과 가설을, III절은 독립된 두 설계를, IV절은 각 사례와 비교 결과를 제시한다.

\section{제도적 배경과 분석틀}
\subsection{스웨덴의 아동보호 행정체계}
2024년 4월 1일 스웨덴 보호시설 제도 개편은 보호자와 동반한 아동을 독립적인 지원 대상으로 다루고 개별적인 필요 판단과 지원 책임을 강화하였다\cite{reform}. 이는 두 사례의 공통 배경이지만, 인구구성 집단 사이에 서로 다른 정책을 배정한 사건은 아니다. 개편을 외생적 충격이나 공통의 처치로 정의하지 않는다.

지방정부의 일반 아동·청소년 조사와 보호시설 동반 아동의 교육보장 업무는 관련 분야에 속하지만 사건 모집단이 다르다. 전자는 일반 사회복지 조사 전체의 처리기간을, 후자는 특정 상황에 필요한 서면 절차를 다룬다. 전국적인 권리 변화와 기관의 보고 상태를 같은 결과로 취급하지 않는다.

\subsection{이주배경 인구구성과 행정 대응의 관계}
인구구성과 행정 대응을 연결할 후보 경로는 언어지원 수요, 기관 사이의 정보 전달, 사건의 위험도와 복잡성, 지역의 업무량 및 전문인력의 축적이다. 조정이 더 필요한 사건이 늘고 인력이 그대로라면 평균 조사기간이 길어질 수 있다. 반대로 숙련과 지원체계가 축적된 지역에서는 인구비중의 증가와 기간의 증가가 함께 나타나지 않을 수 있다. 어느 경로가 작동했는지는 현재의 인구비중과 집계 결과만으로 구분되지 않는다.

부모의 출생지는 국적·체류자격·사용 언어와 동일하지 않다. 지역 주민의 인구구성을 조사를 받은 아동이나 시설 이용자의 속성으로 옮겨 해석해서도 안 된다. 이번 분석의 설명변수는 개인의 처치가 아니라 지역 구성의 지표다.

\subsection{조사기간과 절차 충족의 측정상 차이}
CASE 1에서 기간 증가의 단위는 일수다. CASE 2의 이진 결과에서 집단 격차의 단위는 \%포인트다. 기간이 길어지는 방향과 충족률이 높아지는 방향은 같은 의미를 갖지 않는다. 절차 충족률의 상승은 응답 조건을 충족한 기관이 많아졌다는 뜻이며, 행정 지연이 커졌다는 증거가 아니다. 1에서 0으로의 응답 전환도 실제 절차 폐지인지 조사 기준이나 응답자의 변화인지 구별해야 한다.

\subsection{연구가설}
H1은 이주배경 인구비중과 평균 조사기간 사이에 양의 관계가 있다는 방향 가설이다. H2는 사전 인구구성에 따른 고비중 집단과 비교집단 사이에서 절차 충족률 및 응답 전환이 다르다는 가설이며 방향을 미리 정하지 않는다. H3는 관측 관계가 지역 간 차이, 지역 내 변화, 표본 구성과 측정 방식에 따라 달라질 수 있다는 가설이다. 이 가설들은 관계의 일관성을 검토하기 위한 것으로, 유의한 사양만 골라 채택하지 않는다.

\section{자료 및 실증분석 설계}
\subsection{공통 지역·인구구성 자료와 증거 상태}
두 사례는 네 자리 지방정부 코드를 사용하지만 표본을 통합하지 않는다. CASE 1의 대도시권 분류는 스톡홀름 26개, 예테보리 13개, 말뫼 12개다. CASE 2의 집단 경계는 전국 290개 지방정부의 2015년 해외 출생 비중에서 고정한다. CASE 1의 연도별 아동 배경과 CASE 2의 고정된 사전 집단은 서로 대체할 수 없다.

증거 상태를 구분해 제시한다. CASE 1은 제공 PDF와 대응하는 기존 계산 결과를 유지했다. CASE 2의 표본 수·집단별 충족 건수·최근 전환표는 공식 원자료에서 재구축했고, DML은 새 코드로 독립 재추정했다. 원래 실행 코드와 정제 패널은 확보하지 못했으므로 원래 자료의 모든 행과 폴드까지 같다고 인증하지 않는다. 기존 DR-DID·SDID와 일부 측정·민감도 결과는 원고 보고치로 남겨 두며, 새 추정과 같은 증거로 취급하지 않는다. 아래의 해석은 관측 집계, 새 추정의 방향, 추론의 안정성을 차례로 검토한다.

\subsection{CASE 1 자료와 모형}
\subsubsection{결과변수와 관측 범위}
Kolada U33401은 상반기 0--20세 사회복지 조사의 시작부터 종료까지 평균 일수다. 경제적 지원·가족법·친자확인 조사는 제외된다. 가정폭력 사건이나 이주배경 아동만의 지표가 아니며, 시설 입소 또는 긴급지원 개시까지의 대기시간을 측정하지 않는다. API의 보고 단위 수는 사건 수가 아니므로 사건 가중치로 사용하지 않는다. 51개 지역과 네 연도의 204행 결합틀에서 결측은 그대로 유지하였다.
\begin{table}[htbp]\centering\small
\caption{연도별 평균 조사기간의 관측 범위}\label{tab:coverage}
\begin{tabular}{>{\raggedright\arraybackslash}p{0.16\linewidth} >{\raggedright\arraybackslash}p{0.15\linewidth} >{\raggedright\arraybackslash}p{0.16\linewidth} >{\raggedright\arraybackslash}p{0.17\linewidth} >{\raggedright\arraybackslash}p{0.14\linewidth} >{\raggedright\arraybackslash}p{0.13\linewidth}}
\toprule
연도 & 스톡홀름 & 예테보리 & 말뫼 & 51개 전체 & 전국 \\
\midrule
2022 & 21/26 & 10/13 & 11/12 & 42/51 & 216/290 \\
2023 & 20/26 & 10/13 & 12/12 & 42/51 & 206/290 \\
2024 & 22/26 & 11/13 & 11/12 & 44/51 & 195/290 \\
2025 & 22/26 & 10/13 & 9/12 & 41/51 & 177/290 \\
\bottomrule\end{tabular}
\par\vspace{3pt}\begin{minipage}{.98\linewidth}\footnotesize 각 셀은 관측 지방정부 수/해당 집단 전체 수다. 미관측은 0일이나 조사 미실시를 뜻하지 않는다. 출처: 공식 Kolada 응답의 직접 집계. 제공 CASE 1 원고의 보고값을 유지하였다. 이번 통합에서 재추정하지 않았다.\end{minipage}
\end{table}
2025년의 보고 절차 변경과 연도별 설문 원본의 불완전한 확보는 비교의 한계다. U33401이 유지되었다는 공식 보고와 현장의 측정동등성은 구별한다. 미관측을 조사 미실시나 0일로 대체하지 않는다.

\subsubsection{인구비중과 아동 배경의 분모}
Kolada N02926은 12월 31일 전 연령 전체 주민 중 해외 출생자의 백분율이다. SCB TAB4824는 0--17세 단일 연령과 두 성별을 합산한 A·B·C·D 인원이다. A는 해외 출생 아동, B는 스웨덴 출생·양부모 해외 출생, C는 스웨덴 출생·한 부모만 해외 출생, D는 스웨덴 출생·양부모 스웨덴 출생이다.
\begin{align}
F_{it}&=100\frac{A_{it}+B_{it}}{A_{it}+B_{it}+C_{it}+D_{it}},\\
P_{it}&=100\frac{B_{it}+C_{it}}{B_{it}+C_{it}+D_{it}},\\
Q_{it}&=100\frac{B_{it}}{B_{it}+C_{it}+D_{it}}.
\end{align}
P와 Q의 분모는 스웨덴 출생 아동이므로 A는 제외된다. F의 분모는 전체 아동이다. 공식 스냅샷에 없는 2025년 아동 비중을 생성하지 않았다. 인구의 연말 재고와 상반기 조사기간, 0--17세와 0--20세의 범위 차이도 유지한다. 2015년 비중은 사전 지역구성 비교에, 연도별 값은 관측된 기간의 패널에 각각 사용한다.

\subsubsection{C1-A 횡단면과 C1-B 고정효과}
2024년 회귀는 각 인구지표를 하나씩 넣고 대도시권을 조정한다.
\begin{equation}Y_{i,2024}=\alpha+\beta X_{i,2024}+\gamma_{m(i)}+\varepsilon_i.\end{equation}
HC3 이분산강건 표준오차와 잔차 자유도의 t 구간을 사용한다. 2015년과 2024년의 네 비중은 별도로 추정한다. 패널 모형은
\begin{equation}Y_{it}=\alpha_i+\lambda_t+\beta X_{it}+\varepsilon_{it}\end{equation}
이며 지방정부 군집 표준오차와 원래의 유한표본 보정, t$(G-1)$ 구간을 유지한다. 결과와 설명변수가 함께 한 번만 관측되는 지방정부는 지역 내 기울기에 기여하지 못하므로 제외한다. 시간 불변인 2015년 비중을 지방정부 고정효과와 함께 별도 기울기로 추정하지 않는다. 계수는 비중 1\%포인트당 조사일수다.

\subsubsection{C1-C 차분과 C1-D 영향력}
\begin{equation}\Delta Y_{it}=\delta_t+\beta\Delta X_{it}+\Delta\varepsilon_{it}.\end{equation}
실제로 연결되는 인접 연도 차분만 1년 변화로 사용한다. 2023년 결측을 건너뛴 2022--2024년 변화는 별도의 양 끝 시점 차분이며 HC3 추론을 사용한다. 인접 연도 변화의 결합에는 구간 효과와 지방정부 군집 추론을 사용한다. 균형패널, 대도시권별 연도 효과, 단일 지역과 복수 영향 지역 제외 및 개별 셀 제외를 서로 다른 민감도로 유지한다. 사후 영향력에 따라 고른 제외 표본의 p값은 확인적 검정이 아니다.

\subsection{CASE 2 자료와 모형}
\subsubsection{행정 응답, 집단 경계와 표본}
VIND28의 Ja는 1, Nej는 0, 무응답과 기타 값은 결측이다. 서면 학업보장 절차와 점검 또는 최근 신설 예외의 충족 응답을 다루며 출석·전학 완료·교육 재개일까지의 시간은 포함하지 않는다. 고비중 집단은 2015년 전국 해외 출생 비중의 비가중 75분위 기준값 15.662790386\% 이상으로 정의한다. 이 분류는 정책 처치 배정이 아니다.
\begin{table}[htbp]\centering\small
\caption{CASE 2 전국 기준과 표본 구성}\label{tab:c2flow}
\begin{tabular}{>{\raggedright\arraybackslash}p{0.48\linewidth} >{\raggedright\arraybackslash}p{0.14\linewidth} >{\raggedright\arraybackslash}p{0.14\linewidth} >{\raggedright\arraybackslash}p{0.14\linewidth}}
\toprule
구분 & 고비중 & 비교 & 합계 \\
\midrule
전국 기초자료 & 73 & 217 & 290 \\
비율 단위 두 도시 & 2 & 0 & 2 \\
이진 비교 가능 & 71 & 217 & 288 \\
이력 불완전 & 20 & 77 & 97 \\
장기 균형패널 & 51 & 140 & 191 \\
2024년 공통 표본 & 51 & 139 & 190 \\
2025–2026 연결표본 & 68 & 195 & 263 \\
\bottomrule\end{tabular}
\par\vspace{3pt}\begin{minipage}{.98\linewidth}\footnotesize 집단 경계는 전국 2015년 비중에서 고정. 원고 4–5쪽 및 10–11쪽. 공식 원자료에서 표본·충족 건수·전환 셀을 재구축하여 일치를 확인했다.\end{minipage}
\end{table}
장기패널은 2018--2023년과 2025--2026년의 유효 응답을 같은 191개 시군에서 연결한다. 2024년까지 포함하면 공통 표본은 190개다. 2025--2026년 263개 표본은 최근 두 해의 연결 응답을 위해 별도로 사용한다. 스톡홀름·예테보리는 하위기관 충족 비율을 보고하므로 일반 시군 이진 응답에 합산하지 않는다.

\subsubsection{C2-A 집단 충족률과 C2-B 변화 차이}
집단 지시변수를 $H_i$, 해당 표본 포함을 $S_i$라 하면
\begin{align}
p_{ht}&=E(Y_{it}\mid H_i=h,S_i=1),&G_t&=p_{1t}-p_{0t},\\
\Delta G_{t,s}&=G_t-G_s\nonumber\\
&=E(Y_{it}-Y_{is}\mid H_i=1,S_i=1)-E(Y_{it}-Y_{is}\mid H_i=0,S_i=1).
\end{align}
이는 관측된 두 집단의 변화 차이다. 전국 개혁의 평균처치효과가 아니다. 고비중 집단이 절차 충족에서 우위라는 결과도 더 빠른 사건 처리를 뜻하지 않는다.

\subsubsection{C2-C 응답 전환행렬}
$N_{ab,h}$를 동일 시군의 $a$에서 $b$로의 전환 건수라 하면
\begin{equation}N_h(p_{ht}-p_{hs})=N_{01,h}(s,t)-N_{10,h}(s,t).\end{equation}
0에서 0은 비충족 유지, 0에서 1은 충족 전환, 1에서 0은 비충족 전환, 1에서 1은 충족 유지다. 순변화가 같아도 양방향 전환 건수는 다를 수 있다. 191개와 263개 표본의 포함관계를 이용한 추가 편입 집계 역시 원고에 보고된 값을 유지한다.

\subsubsection{C2-D 조정 방법과 추론}
2015년 기초 공변량은 인구 로그, 노령 비율, 여성 비율, 중위소득과 빈곤 지표다. $e(X_i)=\Pr(H_i=1\mid X_i,S_i=1)$, $m_0(X_i)=E(Y_{it}-Y_{is}\mid H_i=0,X_i,S_i=1)$로 정의한다. $\pi=E(H_i\mid S_i=1)$와 $q_i=(1-H_i)e(X_i)/(1-e(X_i))$를 사용한 정규화 DR 점수는
\begin{equation}\theta_{DR}=E\left[\left\{\frac{H_i}{\pi}-\frac{q_i}{E(q_i\mid S_i=1)}\right\}\{(Y_{it}-Y_{is})-m_0(X_i)\}\mid S_i=1\right].\end{equation}
DR-DID는 두 보조모형을 결합한 조정이다\cite{sz}. DML-DID는 직교 점수와 교차적합을 이용한다\cite{dml}. 원고에서는 21개 주를 5개 폴드로 나누고 평가 주를 학습에서 제외하였다.
\begin{equation}\widehat\theta_{DML}=\frac{1}{n}\sum_{k=1}^{5}\sum_{i\in I_k}\psi(Y_{it}-Y_{is},H_i,X_i;\widehat\eta_{-k}),\quad\widehat\eta_{-k}=(\widehat e_{-k},\widehat m_{0,-k}).\end{equation}
SDID는 과거 응답 경로에 따라 비교 시군과 기간에 가중치를 준다\cite{sdid}. 세 방법은 관측 변화 차이의 조정과 민감도 진단이며, 보조학습이 인과적 정보를 추가로 제공한다고 주장하지 않는다. 특히 DML 수렴률과 학습 안정성은 원고에서 별도로 검증되지 않았다.

비조정 대비는 21개 주를 재표집하는 1,999회 백분위 부트스트랩, DR-DID·DML-DID는 주 군집 영향함수와 자유도 20의 t 구간, SDID는 가중치를 매번 재추정하는 399회 주 부트스트랩을 사용한 보고치다. CASE 1의 지방정부 군집 추론과 이 주 군집 추론은 대체하지 않는다.


\subsubsection{공식 원자료 재구축과 독립 DML-DID 구현}
2026년 11월 18일 Socialstyrelsen의 공식 2026년 Excel DATA 시트와 SCB TAB4823·TAB1121의 2015년 응답을 내려받았다. 다운로드 URL, UTC 시각, 요청 본문, 원본 해시와 정제·추정 코드를 재현 패키지에 보존했다\cite{newraw}. 전국 경계값은 15.662790386208915\%로 일치했고, 장기표본 191개와 연결표본 263개, 여섯 개 집단·연도별 충족 건수와 최근 전환표가 기존 원고와 같았다. 포함·제외 시군의 중위소득 평균도 308.203과 299.400천 SEK로 기존 보고치와 일치했다. 2026년 파일의 역사적 관측을 사용했으며, 당시의 개별 파일 대조와 측정동등성 감사를 새로 완료한 것은 아니다.

Malmö에는 같은 연도의 이진 응답이 반복되고 일부 연도에는 숫자 집계도 함께 있었다. 서로 일치하는 Ja/Nej를 한 건만 채택하고 중복 원행을 남겼다. 상충하는 이진 응답이 나오면 정제를 중단하도록 했다. 다른 도시의 수치 집계를 이진값으로 바꾸거나 결측을 0으로 대체하지 않았다. 원래 코드를 찾지 못했으므로 아래 설정은 이번에 새로 고정한 구현이며, 기존 계수를 맞추기 위해 선택하지 않았다.

공변량은 2015년 인구 로그, 65세 이상 비율, 여성 비율, 소비단위당 중위 가처분소득(자본이득 포함), 저경제수준 비율이다. 소득은 SCB의 2024년 고정가격 천 SEK다. 21개 주를 난수 시드 20261106으로 섞어 5개 평가 폴드로 나누었다. 각 폴드의 주는 두 보조모형의 학습에서 모두 제외했다. 성향점수와 비교집단 변화 회귀에 각각 무작위숲 분류기·회귀기를 사용했으며, 나무 300개, 최소 잎 관측 8개, 전체 공변량 후보를 고정했다. 성향점수는 $[.05,.95]$로 제한했다. 원고의 부트스트랩 시드를 원래 교차적합 시드로 확인한 것은 아니다.

$r_i=\Delta Y_i-\widehat g_{0,-k(i)}(X_i)$, $\widehat q_i=(1-H_i)\widehat e_i/(1-\widehat e_i)$, $\widehat p=\overline H$, $\overline q=n^{-1}\sum_i\widehat q_i$로 쓰면 새 점수는 다음과 같다\cite{dmlscore}.
\begin{align}
\widehat\theta&=\frac1n\sum_i\left(\frac{H_i}{\widehat p}-\frac{\widehat q_i}{\overline q}\right)r_i,\\
\widehat\psi_i&=-\frac{H_i}{\widehat p}\widehat\theta+\left(\frac{H_i}{\widehat p}-\frac{\widehat q_i}{\overline q}\right)r_i.
\end{align}
표준오차는 $\widehat{\mathrm{se}}^2=\{G/(G-1)\}n^{-2}\sum_{c=1}^{G}(\sum_{i\in c}\widehat\psi_i)^2$로 계산하고 $t_{20}$ 구간을 사용했다. 평균 점수 0, 폴드 사이의 주 중복 없음, 공변량 결측 없음, 상수 보조모형에서 단순 DID 점추정·표준오차 일치를 코드로 확인했다. 이는 보조학습 수렴률이나 인과 가정을 검증한 것은 아니다. 패키지의 공개 점수를 Python으로 구현했으며 DoubleML 패키지 자체를 실행한 결과는 아니다.

주 추정 이외에 로지스틱·Ridge와 ExtraTrees, 네 추가 시드, 성향점수 제한 .01·.10, 자본이득 제외 소득을 각각 확인했다. 주 추정과 대안을 합해 기간별 10개 설정이며, 원래 표의 집단 경계 민감도와는 다른 점검이다. 시드 대안은 폴드와 학습기의 난수를 함께 바꾸므로 순수한 폴드 효과로 해석하지 않는다. 세 기간의 변화 보조모형은 각각 학습하므로 독립 추정치에 정확한 가산성을 강제하지 않는다.

\subsection{식별 대상과 통계적 추론의 범위}
단면 계수는 지역 간 선형 관계, 고정효과는 공통 연도 변화와 고정 지역 차이를 제외한 동반 변화, DID는 사전 집단 사이의 응답 변화 차이다. 서로 다른 추정 대상이므로 계수 부호와 통계적 유의성의 차이를 모형 우열로 읽지 않는다. 두 사례의 지역 집계로 개인의 차별적 처리나 개별 보호 지연을 식별할 수 없다. 신뢰구간은 사용한 오차 가정 아래의 불확실성이고, 비무작위 결측과 측정 기준 변경의 불확실성을 모두 포함하는 구간은 아니다.

\FloatBarrier\section{실증분석 결과}
\subsection{CASE 1 인구구성과 사회복지 조사기간}
\subsubsection{기술통계와 횡단면 추정}
2024년 관측 44개 지방정부의 단순 평균은 100.70일, 표준편차는 13.48일이다. 예테보리시 자체의 결과는 결측이므로 대도시권 평균을 그 도시의 성과로 읽지 않는다.
\begin{table}[htbp]\centering\small
\caption{2024년 평균 조사기간의 지역별 분포}\label{tab:descriptives}
\begin{tabular}{>{\raggedright\arraybackslash}p{0.16\linewidth} >{\raggedright\arraybackslash}p{0.15\linewidth} >{\raggedright\arraybackslash}p{0.16\linewidth} >{\raggedright\arraybackslash}p{0.17\linewidth} >{\raggedright\arraybackslash}p{0.14\linewidth} >{\raggedright\arraybackslash}p{0.13\linewidth}}
\toprule
대도시권 & 관측/전체 & 평균 일수 & 표준편차 & 최솟값 & 최댓값 \\
\midrule
스톡홀름 & 22/26 & 103.43 & 13.47 & 86 & 137 \\
예테보리 & 11/13 & 92.00 & 12.80 & 67 & 112 \\
말뫼 & 11/12 & 103.95 & 11.21 & 75 & 117 \\
전체 & 44/51 & 100.70 & 13.48 & 67 & 137 \\
\bottomrule\end{tabular}
\par\vspace{3pt}\begin{minipage}{.98\linewidth}\footnotesize 관측 지방정부별 동일 가중치 통계다. 사건 수 가중 평균이 아니다. 예테보리시의 2024년 값은 미관측이다. 출처: 본고 계산. 제공 CASE 1 원고의 보고값을 유지하였다. 이번 통합에서 재추정하지 않았다.\end{minipage}
\end{table}
네 동시점 인구비중의 원 Pearson 상관은 약 0.015--0.052다. 대도시권을 조정한 네 기울기는 작은 음수이며 신뢰구간은 모두 0을 포함한다. 전체 주민 해외 출생 비중의 계수는 $-0.216$일, 95\% 구간은 $[-0.673,0.241]$일이다. 2015년 비중을 사용한 네 별도 비교도 유의한 양의 관계를 보여 주지 않는다. 8개 기본 단면 검정의 BH 보정값은 모두 약 .419다. 비유의성을 서비스 접근의 동일성으로 해석하지 않는다.
\begin{table}[htbp]\centering\small
\caption{2024년 대도시권 고정효과 단면 회귀}\label{tab:cross}
\begin{tabular}{>{\raggedright\arraybackslash}p{0.4\linewidth} >{\raggedright\arraybackslash}p{0.13\linewidth} >{\raggedright\arraybackslash}p{0.26\linewidth} >{\raggedright\arraybackslash}p{0.1\linewidth}}
\toprule
인구구성 지표 & 계수 & 95\% 구간 & p값 \\
\midrule
전체 주민 해외 출생 & -0.216 & [-0.67, 0.24] & .346 \\
전체 아동 해외배경 & -0.163 & [-0.47, 0.15] & .295 \\
스웨덴 출생 아동 한 명 이상 해외 출생 부모 & -0.123 & [-0.42, 0.17] & .404 \\
스웨덴 출생 아동 양부모 해외 출생 & -0.160 & [-0.48, 0.16] & .325 \\
\bottomrule\end{tabular}
\par\vspace{3pt}\begin{minipage}{.98\linewidth}\footnotesize 각 지표를 별도 추정했다. 모든 모형 n=44, HC3 표준오차와 t(n$-$k) 구간. 단위는 비중 1\%포인트당 일수. 2015년 네 지표를 포함한 기본 단면 8개 검정의 BH 보정값은 모두 약 .419였다. 제공 CASE 1 원고의 보고값을 유지하였다. 이번 통합에서 재추정하지 않았다.\end{minipage}
\end{table}
\subsubsection{지방정부 고정효과와 표본을 맞춘 비교}
패널의 기울기는 모두 양수다. 해외 출생 주민 비중에는 45개 지역·169행, 아동 배경에는 43개 지역·126행이 쓰인다. 전체 아동 해외배경의 기울기는 5.38일, 양부모 해외 출생 비중은 6.70일이다. 원 p값은 각각 .023과 .0048, 네 기본 패널 검정의 BH 보정값은 .046과 .019다.
\begin{table}[htbp]\centering\small
\caption{지방정부와 연도 고정효과의 기본 패널 추정}\label{tab:panel}
\begin{tabular}{>{\raggedright\arraybackslash}p{0.35\linewidth} >{\raggedright\arraybackslash}p{0.12\linewidth} >{\raggedright\arraybackslash}p{0.12\linewidth} >{\raggedright\arraybackslash}p{0.25\linewidth} >{\raggedright\arraybackslash}p{0.09\linewidth}}
\toprule
인구구성 지표 & n/G & 계수 & 95\% 구간 & p값 \\
\midrule
전체 주민 해외 출생 & 169/45 & 5.78 & [-0.92, 12.49] & .089 \\
전체 아동 해외배경 & 126/43 & 5.38 & [0.78, 9.97] & .023 \\
스웨덴 출생 아동 한 명 이상 해외 출생 부모 & 126/43 & 4.60 & [-0.94, 10.14] & .101 \\
스웨덴 출생 아동 양부모 해외 출생 & 126/43 & 6.70 & [2.16, 11.23] & .005 \\
\bottomrule\end{tabular}
\par\vspace{3pt}\begin{minipage}{.98\linewidth}\footnotesize G는 지방정부 수. 첫 행은 2022–2025년, 나머지는 2022–2024년이다. 지방정부 군집 표준오차, t(G$-$1) 구간. 단위는 1\%포인트당 일수. 해외배경·양부모 해외 출생의 BH 보정 p값은 각각 .046, .019다. 제공 CASE 1 원고의 보고값을 유지하였다. 이번 통합에서 재추정하지 않았다.\end{minipage}
\end{table}
동일 40개 지방정부로 맞추어도 양부모 비중의 단면 계수는 $-0.105$일, 패널 계수는 $+6.555$일이다. 표본의 지역 구성만으로 부호 차이를 설명할 수는 없다. 다만 이 표본에서 지역·연도 효과 제거 후 비중의 잔차 표준편차는 약 0.27\%포인트로 작다. 시간가변 인력이나 사례 구성이 무엇을 설명하는지는 구분되지 않는다.
\begin{table}[htbp]\centering\small
\caption{양부모 해외 출생 비중의 패널 민감도}\label{tab:panel_sensitivity}
\begin{tabular}{>{\raggedright\arraybackslash}p{0.35\linewidth} >{\raggedright\arraybackslash}p{0.12\linewidth} >{\raggedright\arraybackslash}p{0.12\linewidth} >{\raggedright\arraybackslash}p{0.25\linewidth} >{\raggedright\arraybackslash}p{0.09\linewidth}}
\toprule
모형 & n/G & 계수 & 95\% 구간 & p값 \\
\midrule
기본 불균형패널 & 126/43 & 6.70 & [2.16, 11.23] & .005 \\
균형패널 & 120/40 & 6.55 & [0.37, 12.74] & .038 \\
대도시권별 연도 효과 & 126/43 & 6.52 & [0.10, 12.94] & .047 \\
균형표본과 대도시권별 연도 효과 & 120/40 & 6.07 & [-2.09, 14.24] & .140 \\
\bottomrule\end{tabular}
\par\vspace{3pt}\begin{minipage}{.98\linewidth}\footnotesize 모든 행은 2022–2024년이다. 지방정부 군집 표준오차와 유한표본 보정을 적용했다. 민감도 p값은 탐색적이며 다중검정 보정 전 값이다. 제공 CASE 1 원고의 보고값을 유지하였다. 이번 통합에서 재추정하지 않았다.\end{minipage}
\end{table}
균형표본과 대도시권별 연도 효과를 함께 사용하면 계수는 6.07일, 구간은 $[-2.09,14.24]$일, p값은 .140이다. 양의 점추정치의 유지와 크기의 정밀성은 구별된다.
\subsubsection{일차차분과 영향력 진단}
43개 지방정부를 하나씩 제외한 기울기는 모두 양수이며 5.64--7.17일이다. 42회는 원 p값이 .05 미만이지만 Upplands-Bro 제외에서는 p=.057이다. Öckerö·Nynäshamn·Upplands-Bro를 함께 제외하면 3.21일, p=.263으로 낮아진다. 세 지역의 조정 후 부호 있는 교차곱 기여 71.1\%는 사건 수 비중이나 설명된 분산의 비율이 아니다.
\begin{table}[htbp]\centering\small
\caption{양부모 해외 출생 비중의 지방정부 제외 진단}\label{tab:influence}
\begin{tabular}{>{\raggedright\arraybackslash}p{0.21\linewidth} >{\raggedright\arraybackslash}p{0.25\linewidth} >{\raggedright\arraybackslash}p{0.43\linewidth}}
\toprule
제외한 지방정부 & 재추정 계수 & p값 \\
\midrule
제외 없음 & 6.70 & .0048 \\
Öckerö 1407 & 5.64 & .007 \\
Nynäshamn 0192 & 6.03 & .017 \\
Upplands-Bro 0139 & 6.04 & .057 \\
Staffanstorp 1230 & 6.11 & .012 \\
Burlöv 1231 & 6.37 & .008 \\
Öckerö·Nynäshamn·Upplands-Bro 동시 제외 & 3.21 & .263 \\
상위 5개 지지 지방정부 동시 제외 & 1.10 & .690 \\
\bottomrule\end{tabular}
\par\vspace{3pt}\begin{minipage}{.98\linewidth}\footnotesize 기본 지역·연도 고정효과 모형을 매번 재추정했다. 집단 제외 순위는 전체 표본의 단일 제외에서 계수를 가장 낮추는 순서로 정한 사후 스트레스 테스트다. 상위 5곳은 앞의 세 곳에 Staffanstorp와 Burlöv를 더한 집단이다. 제공 CASE 1 원고의 보고값을 유지하였다. 이번 통합에서 재추정하지 않았다.\end{minipage}
\end{table}
Öckerö 2022년과 Tjörn 2024년 셀을 함께 제외하면 4.81일(p=.069), Nynäshamn 2022년까지 제외하면 3.99일(p=.164)이다. Upplands-Bro는 두 시점만 있어 한 셀 제외가 남은 단일 관측의 제외도 유발한다. 영향력이 크다는 사실로 공식 관측값을 오류로 판정하지 않는다.
\begin{table}[htbp]\centering\small
\caption{양부모 해외 출생 비중의 일차차분 추정}\label{tab:fd}
\begin{tabular}{>{\raggedright\arraybackslash}p{0.35\linewidth} >{\raggedright\arraybackslash}p{0.12\linewidth} >{\raggedright\arraybackslash}p{0.12\linewidth} >{\raggedright\arraybackslash}p{0.25\linewidth} >{\raggedright\arraybackslash}p{0.09\linewidth}}
\toprule
모형 & 변화/G & 계수 & 95\% 구간 & p값 \\
\midrule
인접 연도 차분 & 82/42 & 4.25 & [-2.83, 11.33] & .232 \\
인접 연도와 대도시권별 구간 효과 & 82/42 & 4.06 & [-4.32, 12.45] & .333 \\
2022–2024년 차분 & 41/41 & 7.74 & [4.42, 11.07] & <.001 \\
2022–2024년 차분과 대도시권 조정 & 41/41 & 7.43 & [2.31, 12.55] & .006 \\
\bottomrule\end{tabular}
\par\vspace{3pt}\begin{minipage}{.98\linewidth}\footnotesize 연간 차분에는 실제 연속 관측만 사용하며 구간 효과를 포함한다. 연간 차분은 지방정부 군집, 양 끝 시점 차분은 HC3 표준오차를 사용했다. 비중 1\%포인트당 일수이며 두 시점 차분을 1년 변화로 간주하지 않았다. 제공 CASE 1 원고의 보고값을 유지하였다. 이번 통합에서 재추정하지 않았다.\end{minipage}
\end{table}
연간 차분은 42개 지역의 82개 변화에서 4.25일, 구간 $[-2.83,11.33]$일, p=.232다. 반면 2022--2024년 양 끝 시점 차분은 7.74일이고 대도시권을 조정하면 7.43일, 구간 $[2.31,12.55]$일, p=.0056이다. 후자는 41개 지역을 각각 제외해도 양수와 원 p값 .05 미만을 유지하지만, Staffanstorp·Österåker·Nynäshamn을 함께 제외하면 4.64일(p=.082)이다. 균형표본의 연간 차분도 정밀하지 않아 표본 차이만으로 결과 차이를 설명할 수 없다.
\subsubsection{CASE 1 결과의 요약}
H1을 지역 전반의 규칙으로 받아들일 근거는 부족하다. 2024년 지역 간 기울기는 작지만, 같은 지역의 일부 인구비중 변화에서는 평균 조사기간과 양의 관계가 관측된다. 양부모 해외 출생 비중의 관계는 한 지역만의 제외로 사라지지는 않으나, 복수 영향 지역을 제외하거나 연간 차분을 사용하면 불확실성이 커진다. 따라서 특정 패널 모형에서의 양의 연관성은 보고할 수 있지만, 이주배경 비중 증가가 조사기간을 늘렸다고 해석할 수는 없다. 단면의 비유의성 역시 기간 차이가 없다는 증명은 아니다.

\FloatBarrier\subsection{CASE 2 인구 집중도와 학업보장 절차}
\subsubsection{연도별 충족률과 집단 간 격차}
\begin{table}[htbp]\centering\small
\caption{CASE 2 장기패널의 집단 충족 건수와 격차}\label{tab:c2levels}
\begin{tabular}{>{\raggedright\arraybackslash}p{0.13\linewidth} >{\raggedright\arraybackslash}p{0.24\linewidth} >{\raggedright\arraybackslash}p{0.25\linewidth} >{\raggedright\arraybackslash}p{0.25\linewidth}}
\toprule
연도 & 고비중 (51) & 비교 (140) & 격차 (\%p) \\
\midrule
2023 & 17/51 & 39/140 & +5.48 \\
2025 & 31/51 & 56/140 & +20.78 \\
2026 & 27/51 & 69/140 & +3.66 \\
\bottomrule\end{tabular}
\par\vspace{3pt}\begin{minipage}{.98\linewidth}\footnotesize 인원은 원고에 보고된 동일 표본 충족 건수. 격차는 고비중–비교집단. 공식 원자료에서 표본·충족 건수·전환 셀을 재구축하여 일치를 확인했다.\end{minipage}
\end{table}
고비중 집단의 충족 건수는 17에서 31로 늘어난 뒤 27로 줄고, 비교집단은 39에서 56, 69로 증가한다. 따라서 격차는 2025년 확대되었다가 2026년에 좁혀진다. 2026년 두 집단의 절차가 동일하다거나 이후 수렴이 지속된다는 의미는 아니다.
\subsubsection{DID와 조정 추정}
\begin{table}[htbp]\centering\small
\caption{CASE 2 기존 원고의 보고치: 변화 차이와 조정 추정}\label{tab:c2did}
\begin{tabular}{>{\raggedright\arraybackslash}p{0.22\linewidth} >{\raggedright\arraybackslash}p{0.21\linewidth} >{\raggedright\arraybackslash}p{0.16\linewidth} >{\raggedright\arraybackslash}p{0.16\linewidth} >{\raggedright\arraybackslash}p{0.16\linewidth}}
\toprule
기간 & 방법 & 추정치 & 하한 & 상한 \\
\midrule
2025–2023 & 비조정 & +15.31 & +2.37 & +27.97 \\
 & DR-DID & +14.31 & $-$2.50 & +31.12 \\
 & DML-DID & +18.10 & +1.81 & +34.39 \\
 & SDID & +16.27 & +4.54 & +27.37 \\
2026–2023 & 비조정 & $-$1.82 & $-$17.26 & +16.64 \\
 & DR-DID & $-$1.32 & $-$28.72 & +26.09 \\
 & DML-DID & $-$8.36 & $-$25.68 & +8.96 \\
 & SDID & +0.54 & $-$12.77 & +14.47 \\
2026–2025 & 비조정 & $-$17.13 & $-$36.82 & +0.37 \\
 & DR-DID & $-$15.62 & $-$37.21 & +5.96 \\
 & DML-DID & $-$26.46 & $-$44.38 & $-$8.54 \\
 & SDID & $-$15.73 & $-$34.54 & +3.77 \\
\bottomrule\end{tabular}
\par\vspace{3pt}\begin{minipage}{.98\linewidth}\footnotesize 단위 \%p, n=191, 주 21개. 비조정 1,999회 주 부트스트랩; DR·DML 주 영향함수 t(20); SDID 399회 주 부트스트랩. 원고 표 5. 제공 CASE 2 원고 보고치의 전사이며 이번 통합에서 재추정하지 않았다.\end{minipage}
\end{table}
기존 원고의 DML-DID는 2025--2023년 $+18.10$, 2026--2025년 $-26.46$\%포인트이고 두 보고 구간은 0을 제외했다. 그러나 이 수치를 공식 원자료의 새 구현으로 그대로 얻지는 못했다. 표~\ref{tab:c2new}는 원래 코드의 정확 재현을 주장하지 않고, 원고와 표본 수·구성요건 및 기초 공변량을 맞춘 독립 재추정과 나란히 비교한다.

\begin{table}[htbp]\centering\small
\caption{CASE 2 기존 DML 보고치와 공식 원자료의 독립 재추정}\label{tab:c2new}
\begin{tabular}{>{\raggedright\arraybackslash}p{0.21\linewidth} >{\raggedright\arraybackslash}p{0.14\linewidth} >{\raggedright\arraybackslash}p{0.15\linewidth} >{\raggedright\arraybackslash}p{0.29\linewidth} >{\raggedright\arraybackslash}p{0.1\linewidth}}
\toprule
대비 & 기존 보고 & 독립 추정 & 새 95\% 구간 & p값 \\
\midrule
2025-2023 & +18.10 & +9.24 & [-16.18, 34.67] & 0.457 \\
2026-2023 & -8.36 & -11.21 & [-35.83, 13.41] & 0.354 \\
2026-2025 & -26.46 & -18.43 & [-40.09, 3.22] & 0.091 \\
\bottomrule\end{tabular}
\par\vspace{3pt}\begin{minipage}{.98\linewidth}\footnotesize 단위 \%p; n=191, 고비중 51, 주 G=21. 새 구간은 주 영향함수 표준오차와 t(20). 기존 코드를 발견하지 못했으므로 정확 일치 재현이 아니다. 기존 구간은 앞 표의 보고치와 구분한다.\end{minipage}
\end{table}\begin{table}[htbp]\centering\small
\caption{CASE 2 독립 DML 구현의 민감도}\label{tab:c2newrobust}
\begin{tabular}{>{\raggedright\arraybackslash}p{0.21\linewidth} >{\raggedright\arraybackslash}p{0.29\linewidth} >{\raggedright\arraybackslash}p{0.18\linewidth} >{\raggedright\arraybackslash}p{0.23\linewidth}}
\toprule
대비 & 계수 범위 (\%p) & 양수/음수 & p<.05 / 설정 수 \\
\midrule
2025-2023 & [9.21, 15.07] & 10/0 & 0/10 \\
2026-2023 & [-12.78, 1.34] & 1/9 & 0/10 \\
2026-2025 & [-23.77, -13.73] & 0/10 & 1/10 \\
\bottomrule\end{tabular}
\par\vspace{3pt}\begin{minipage}{.98\linewidth}\footnotesize 기간별 주 추정과 대안 9개. 대안은 학습기 2개, 시드 4개, 절단 2개, 소득 개념 1개다. 서로 독립인 검정이 아니며 다중검정 보정 전 값이다. 기존 원고의 민감도 표를 재현한 결과가 아니다.\end{minipage}
\end{table}

공식 자료의 독립 재추정에서는 2025--2023년 $+9.24$\%포인트, 95\% 구간 $[-16.18,34.67]$이었고, 2026--2023년은 $-11.21$ $[-35.83,13.41]$, 2026--2025년은 $-18.43$ $[-40.09,3.22]$였다. 세 주 구간이 모두 0을 포함하므로, 원래의 두 유의한 DML 대비가 이 주 분할에서도 그대로 확인됐다고 기술할 수 없다. 원래 학습기와 폴드가 불명확하여 차이의 원인을 특정 설정 하나로 귀속하거나 기존 수치가 오류라고 판정할 수도 없다.

가중치 집중과 분할 변동을 살펴볼 이유도 있다. 새 비교집단 유효표본크기는 54.79였으며 성향점수 21개가 .05 아래에서 제한되었다. 기간별 10개 설정 모두에서 2025--2023년은 양수, 2026--2025년은 음수였지만 전자는 유의한 설정이 없고 후자는 한 설정만 $p<.05$였다. 2026--2023년은 학습기에 따라 부호도 바뀌었다. 이 첫 점검은 방향과 유의성을 구분할 필요를 보여 준다. 다만 시드 변경이 폴드와 학습기 난수를 함께 바꾸므로, 어느 요인이 변동을 만드는지는 다음의 분리 점검으로 확인한다.


\begin{table}[htbp]\centering\small
\caption{CASE 2 추가 점검: 폴드 배정과 학습기 난수의 분리}\label{tab:c2second}
\begin{tabular}{>{\raggedright\arraybackslash}p{0.24\linewidth} >{\raggedright\arraybackslash}p{0.19\linewidth} >{\raggedright\arraybackslash}p{0.29\linewidth} >{\raggedright\arraybackslash}p{0.19\linewidth}}
\toprule
변경 요인 & 대비 & 계수 범위 (\%p) & p<.05 / 횟수 \\
\midrule
폴드만 변경 & 2025-2023 & [5.79, 24.92] & 2/10 \\
폴드만 변경 & 2026-2023 & [-24.14, 1.12] & 0/10 \\
폴드만 변경 & 2026-2025 & [-31.50, -17.96] & 8/10 \\
학습 난수만 변경 & 2025-2023 & [10.00, 10.93] & 0/5 \\
학습 난수만 변경 & 2026-2023 & [-11.61, -8.83] & 0/5 \\
학습 난수만 변경 & 2026-2025 & [-21.01, -18.70] & 1/5 \\
\bottomrule\end{tabular}
\par\vspace{3pt}\begin{minipage}{.98\linewidth}\footnotesize 표본 191개·주 21개와 학습기 설정을 유지. 폴드만 변경은 학습 난수를 고정한 10개 배정, 학습 난수만 변경은 폴드를 고정한 5개 난수다. 구간은 주 군집 t(20). 빈도는 보정 전이며 서로 독립인 검정이나 통합 유의확률이 아니다. 원래 폴드를 복원한 결과도 아니다.\end{minipage}
\end{table}

추가로 폴드 배정과 학습기 자체의 난수를 분리하여 검사했다. 동일 폴드에서 성향점수 모형만 로지스틱으로 바꾸면 2025--2023년 대비는 9.24에서 13.85\%포인트, 결과 변화 모형만 Ridge로 바꾸면 13.01\%포인트가 되었다. 두 모형을 함께 바꾸면 15.07\%포인트였다. 반면 고정 예측에서 비교 가중치의 정규화만 바꾼 진단은 9.59\%포인트로 변화가 작았다. 두 표준오차 진단에서도 기본 추정의 유의성 판정은 달라지지 않았다. 이 비교는 새 구현 내부의 민감도를 분리하며, 불명확한 원래 구현과의 차이를 특정 원인으로 확정하지 않는다.

특히 주별 폴드 배정만 바꾼 10개 설정에서 2026--2025년 계수는 모두 음수였으며 8개 구간이 0을 제외했다. 2025--2023년은 모두 양수였지만 유의한 설정은 2개였고, 2026--2023년은 어느 설정에서도 유의하지 않았다. 고정 폴드에서 학습기 난수만 바꿀 때보다 계수 변동이 컸다. 따라서 한 분할에서의 비유의를 신호의 부재로 일반화해서도 안 되고, 유의한 분할만 골라 결과를 확정해서도 안 된다. 단기 격차 축소의 방향은 반복되며 그 크기와 유의성은 분할에 의존한다. 8/10은 귀무가설이 틀릴 확률이나 새로운 통합 p값이 아니다. 검토한 10개 분할을 가능한 모든 분할의 대표 표본으로 보장하지 않으며 계수의 최소·최대도 신뢰구간이 아니다. 2026--2023년의 비유의성도 변화가 정확히 0이거나 2023년 상태로 완전히 복귀했다는 검정 결과가 아니다. 이 점검은 사후 민감도 진단이며 원래 분석의 정확 재현을 대신하지 않는다.


원고의 비교집단 유효표본크기는 DR-DID 14.79, DML-DID 84.60, 기존 평균 SDID 138.81로 보고된다. 가중 규칙이 서로 달라 구간 폭만으로 우열을 정할 수 없다. 비조정 대비의 주 하나씩 제외한 점추정 범위는 2025--2023년 $[12.44,17.72]$, 2026--2023년 $[-6.03,1.60]$, 2026--2025년 $[-20.46,-10.84]$\%포인트이며 이는 신뢰구간이 아니다.

\subsubsection{응답 전환행렬}
\begin{table}[htbp]\centering\small
\caption{CASE 2 연결 응답 전환}\label{tab:c2transitions}
\begin{tabular}{lrrrrrr}
\toprule
기간 & 집단 & 0$\to$0 & 0$\to$1 & 1$\to$0 & 1$\to$1 & 순변화 \\
\midrule
2023$\to$2025 & 비교 & 68 & 33 & 16 & 23 & +17 \\
2023$\to$2025 & 고비중 & 16 & 18 & 4 & 13 & +14 \\
2025$\to$2026 & 비교 & 51 & 33 & 20 & 36 & +13 \\
2025$\to$2026 & 고비중 & 14 & 6 & 10 & 21 & $-$4 \\
\bottomrule\end{tabular}
\par\vspace{3pt}\begin{minipage}{.98\linewidth}\footnotesize 원고 표 6. 첫 기간은 2년, 둘째는 1년. 코드 전환을 실제 절차의 신설·폐지로 단정하지 않음. 공식 원자료에서 표본·충족 건수·전환 셀을 재구축하여 일치를 확인했다.\end{minipage}
\end{table}
2025--2026년 고비중 집단의 순감소 4건은 충족 전환 6건과 비충족 전환 10건의 차이다. 비교집단 순증가 13건도 33건과 20건의 이동을 포함한다. 두 집단을 합친 69개 시군, 191개 중 36.1\%가 코드를 바꾸었다. 2023--2025년 전환은 두 해 사이의 이동으로, 1년 간격의 전환 건수와 연간 발생률처럼 비교하지 않는다.

\subsubsection{표본 확장과 측정 민감도}
263개 연결표본은 최근 두 해를 비교하는 데 장기 완전 응답 요건을 부과했을 때 무엇이 달라지는지 보여 준다. 표본 확장 뒤에도 비조정 변화 차이는 $-11.79$\%포인트로 격차 축소 방향을 유지한다. 그러나 고비중 집단의 충족 건수는 35개에서 35개로 같고, 비교집단은 70개에서 93개로 늘어난다. 단기 격차 축소와 고비중 집단의 절차 후퇴는 서로 다른 주장이다. 후자는 장기표본의 응답 감소만으로 전체 연결표본에 일반화할 수 없다.

\begin{table}[htbp]\centering\small
\caption{CASE 2 표본 확장에 따른 응답 변화}\label{tab:c2expand}
\begin{tabular}{lrrrrrrr}
\toprule
집단 & 표본 & n & 2025년 1 & 2026년 1 & 순변화 & 0$\to$1 & 1$\to$0 \\
\midrule
고비중 & 장기 & 51 & 31 & 27 & $-$4 & 6 & 10 \\
 & 추가 & 17 & 4 & 8 & +4 & 7 & 3 \\
 & 연결 & 68 & 35 & 35 & 0 & 13 & 13 \\
비교 & 장기 & 140 & 56 & 69 & +13 & 33 & 20 \\
 & 추가 & 55 & 14 & 24 & +10 & 15 & 5 \\
 & 연결 & 195 & 70 & 93 & +23 & 48 & 25 \\
\bottomrule\end{tabular}
\par\vspace{3pt}\begin{minipage}{.98\linewidth}\footnotesize 원고 표 7의 보고값. 추가 편입 건수도 원고의 산술 분해이며 이번에 새로운 전환 자료를 만들지 않음. 공식 원자료에서 표본·충족 건수·전환 셀을 재구축하여 일치를 확인했다.\end{minipage}
\end{table}
263개 연결표본의 고비중 집단은 35에서 35로 유지된다. 추가 편입 17개 시군의 4건 증가가 장기패널의 4건 감소를 상쇄한다. 비교집단에는 추가 55개 시군의 10건 증가가 더해진다. 확대 표본의 격차 변화는 $-11.79$\%포인트, 주 군집 구간은 $[-26.47,6.24]$\%포인트다. 표본 확장은 단순히 정밀성만 바꾸지 않고 관측 변화의 구성을 바꾼다.
\begin{table}[htbp]\centering\small
\caption{CASE 2 측정동등성 감사의 보고 범위}\label{tab:c2audit}
\begin{tabular}{>{\raggedright\arraybackslash}p{0.28\linewidth} >{\raggedright\arraybackslash}p{0.64\linewidth}}
\toprule
항목 & 제공 원고의 보고 내용 \\
\midrule
문항과 결정 수준 & 서면 학업보장 절차와 ledningsnivå 조건 확인 \\
점검 기간 & 24개월 표기, 12개월 변경 증거 없음 \\
최근 신설 예외 & 존재 확인, 연도별 기준일 동등성 미확인 \\
역사적 공개값 & 2023–2025 공통 유효 값 일치 보고 \\
과거 조사표·기관 경계 & 2018–2022 개별 인증과 전 기간 조직 감사 미완료 \\
\bottomrule\end{tabular}
\par\vspace{3pt}\begin{minipage}{.98\linewidth}\footnotesize 원고 표 3 및 본문. 전체 시계열의 완전 측정동등성은 미확인(U). 제공 CASE 2 원고 보고치의 전사이며 이번 통합에서 재추정하지 않았다.\end{minipage}
\end{table}
원고는 2023년 독립 결과표의 유효 273개 시군과 2026년 파일의 같은 연도 값이 일치하고, 2024·2025년 공통 유효 키에서도 불일치가 없었다고 보고한다. 그러나 최근 신설 예외의 기준일과 2018--2022년 당시 조사표는 완전히 인증되지 않았다. 동일 코드의 일치와 전 기간 운영상태의 측정동등성은 다르다. 이 통합에서는 감사 원파일이 없으므로 그 점검 자체를 새로 재현했다고 주장하지 않는다.
\begin{table}[htbp]\centering\small
\caption{CASE 2 결측 응답의 논리적 범위}\label{tab:c2missing}
\begin{tabular}{>{\raggedright\arraybackslash}p{0.12\linewidth} >{\raggedright\arraybackslash}p{0.25\linewidth} >{\raggedright\arraybackslash}p{0.25\linewidth} >{\raggedright\arraybackslash}p{0.28\linewidth}}
\toprule
연도 & 비교 비율 (\%) & 고비중 비율 (\%) & 격차 (\%p) \\
\midrule
2023 & 22.58–27.19 & 32.39–39.44 & +5.21–+16.86 \\
2024 & 27.65–32.72 & 36.62–36.62 & +3.90–+8.97 \\
2025 & 33.18–38.71 & 50.70–53.52 & +11.99–+20.34 \\
2026 & 45.62–51.61 & 49.30–50.70 & $-$2.32–+5.08 \\
\bottomrule\end{tabular}
\par\vspace{3pt}\begin{minipage}{.98\linewidth}\footnotesize 이진 비교 가능 288개 기준. 원고 표 8. 신뢰구간과 구별. 제공 CASE 2 원고 보고치의 전사이며 이번 통합에서 재추정하지 않았다.\end{minipage}
\end{table}
결측 범위는 미응답이 모두 비충족 또는 충족일 수 있는 논리적 경계다.
\begin{equation}\frac{A_{ht}}{N_h}\leq p^{all}_{ht}\leq\frac{A_{ht}+M_{ht}}{N_h}.\end{equation}
여기서 $A$는 유효 충족 건수이고 $M$은 결측 건수로, CASE 1의 아동 배경 A와 다른 기호다. 구간은 표본추출 신뢰구간이 아니고 응답을 실제로 대체한 새 데이터도 아니다.
\begin{table}[htbp]\centering\small
\caption{CASE 2 동일 190개 시군의 2024년 처리 민감도}\label{tab:c2year}
\begin{tabular}{>{\raggedright\arraybackslash}p{0.26\linewidth} >{\raggedright\arraybackslash}p{0.22\linewidth} >{\raggedright\arraybackslash}p{0.22\linewidth} >{\raggedright\arraybackslash}p{0.22\linewidth}}
\toprule
방법 & 2023 기준 & 2024 기준 & 변화 \\
\midrule
비조정 & +6.62 & +7.93 & +1.31 \\
DR-DID & +6.49 & +16.07 & +9.58 \\
DML-DID & +4.48 & +8.16 & +3.68 \\
SDID & +8.07 & +8.03 & $-$0.04 \\
\bottomrule\end{tabular}
\par\vspace{3pt}\begin{minipage}{.98\linewidth}\footnotesize 단위 \%p, 2025–2026 사후 평균 대비. 원고 표 9. DR·DML은 기준 응답, SDID는 과거 적합 기간 변경. 제공 CASE 2 원고 보고치의 전사이며 이번 통합에서 재추정하지 않았다.\end{minipage}
\end{table}
같은 190개 시군에서 기준을 바꾸면 DR-DID는 6.49에서 16.07, DML-DID는 4.48에서 8.16\%포인트로 달라진다. SDID는 8.07과 8.03으로 가깝다. DR·DML은 기준 응답을 바꾸고 SDID는 과거 적합 기간을 바꾸므로 동일 효과의 추정 오차라고 보지 않는다.
\begin{table}[htbp]\centering\small
\caption{CASE 2 집단 경계별 사후 평균 대비}\label{tab:c2threshold}
\begin{tabular}{lrrrrrr}
\toprule
분위 & 경계 (\%) & 고비중 n & 비조정 & DR-DID & DML-DID & SDID \\
\midrule
50 & 11.776 & 102 & $-$5.17 & +10.53 & +0.75 & $-$3.48 \\
60 & 13.368 & 83 & $-$1.97 & +8.56 & +2.51 & $-$0.90 \\
70 & 14.976 & 60 & +6.92 & +16.01 & +11.08 & +7.73 \\
75 & 15.663 & 51 & +6.74 & +6.50 & +4.70 & +8.40 \\
80 & 16.589 & 41 & +8.94 & +7.31 & +3.33 & +6.57 \\
\bottomrule\end{tabular}
\par\vspace{3pt}\begin{minipage}{.98\linewidth}\footnotesize 경계 이외 계수 단위 \%p. 동일 191개에서 매번 집단 재분류. 원고 표 10. 제공 CASE 2 원고 보고치의 전사이며 이번 통합에서 재추정하지 않았다.\end{minipage}
\end{table}
연속 해외 출생 비중 10\%포인트와 관련된 조정 회귀의 2025년 대비는 $-0.14$\%포인트, 구간 $[-15.23,14.95]$이며 2026년은 $-2.16$, 구간 $[-17.23,12.91]$이다. 집단 경계와 연속 기울기는 다른 추정 대상이다.
\begin{table}[htbp]\centering\small
\caption{CASE 2 대도시의 서로 다른 공개 단위}\label{tab:c2urban}
\begin{tabular}{>{\raggedright\arraybackslash}p{0.18\linewidth} >{\raggedright\arraybackslash}p{0.12\linewidth} >{\raggedright\arraybackslash}p{0.12\linewidth} >{\raggedright\arraybackslash}p{0.12\linewidth} >{\raggedright\arraybackslash}p{0.36\linewidth}}
\toprule
도시 & 2023 & 2025 & 2026 & 단위 \\
\midrule
스톡홀름 & 50.00 & 50.00 & 72.73 & 하위기관 충족 비율 (\%) \\
말뫼 & 0 & 1 & 1 & 시군 이진값 \\
예테보리 & 75.00 & 100.00 & 0.00 & 하위기관 충족 비율 (\%) \\
\bottomrule\end{tabular}
\par\vspace{3pt}\begin{minipage}{.98\linewidth}\footnotesize 원고 표 12. 단위를 합산하거나 임의로 이진화하지 않음. 제공 CASE 2 원고 보고치의 전사이며 이번 통합에서 재추정하지 않았다.\end{minipage}
\end{table}
기본 CASE 2에서 스톡홀름·예테보리는 이미 비율 단위로 분리되어 있다. 세 도시를 제외한 190개 장기패널의 2026--2023년 DR-DID·DML-DID·SDID 대비는 각각 $-3.15$, $-10.29$, $-0.88$\%포인트로 보고되며 세 주 군집 구간은 모두 0을 포함한다. 공개 비율을 임의로 0과 1로 바꾸거나 일반 시군 값과 평균하지 않는다.

\subsubsection{CASE 2 결과의 요약}
H2와 관련해 직접 확인되는 것은 집단별 공개 응답의 수준과 시간 경로 차이다. 고비중 집단의 충족률은 세 기준 시점 모두 더 높았고, 집단 간 격차는 2025년까지 확대된 뒤 2026년에 축소됐다. 새 DML의 2026--2025년 주 구간은 0을 포함하지만, 폴드만 바꾼 10개 설정에서는 모두 음수이고 8개가 유의했다. 이는 단기 축소의 방향에 반복적인 지지가 있음을 보여 주며, 그 크기나 유의성이 분할과 무관하게 확정됐다는 뜻은 아니다. 같은 점검에서 2026--2023년은 부호가 달라지고 유의한 설정이 없어 지속적인 변화의 방향을 안정적으로 판단하기 어렵다.

263개 연결표본에서도 격차 축소는 관측되지만 고비중 집단의 충족 건수 감소는 나타나지 않는다. 그러므로 단기 격차의 축소는 두 표본에서 관측되는 결과이고, 고비중 집단의 일방적인 절차 후퇴는 표본에 따라 달라지는 해석이다. 어느 결과도 절차가 실제로 폐지됐거나 아동의 지원이 지연됐음을 직접 확인하지 않는다.

\FloatBarrier\subsection{두 사례의 실증 결과 비교}
\subsubsection{처리기간과 절차 충족의 방향}
CASE 1의 양의 계수는 평균 조사기간과 인구비중의 양의 관계이고, CASE 2의 양의 격차는 고비중 집단의 절차 충족 응답이 더 높다는 뜻이다. 두 양수는 같은 방향의 성과나 지연을 나타내지 않는다. CASE 2의 2025--2026년 격차 축소를 CASE 1의 일수 증가에 더하거나 동일 결과로 평균할 수 없다.

\subsubsection{지역 간 차이와 시간에 따른 변화}
CASE 1에서는 단면과 동일 지역 변화의 기울기가 달랐고, CASE 2에서는 2025년과 2026년의 집단 격차 및 응답 전환이 달랐다. 시간 변화를 사용했다는 이유만으로 두 추정이 같아지지는 않는다. CASE 1은 연도별 연속 인구비중, CASE 2는 2015년 고정 집단으로 분류하며, 각각 지방정부 군집과 주 군집 추론을 사용한다.

\subsubsection{표본과 측정에 따른 안정성}
안정성은 세 층위에서 다르게 나타난다. CASE 2의 집단별 충족 건수와 전환은 공식 자료에서 확인된 관측 사실이다. 2025--2026년 격차 축소는 두 표본에서 관측되고 추가 폴드 점검에서도 방향이 유지된다. 반면 보정된 계수의 크기와 유의성은 설정에 따라 달라진다. 세 층위를 한꺼번에 강건하거나 불안정하다고 묶으면, 반복되는 방향과 불확실한 추론을 모두 놓치게 된다.

CASE 1의 복수 영향 지역 제외는 패널 기울기를 지지하는 정보의 집중을 보여 주고, CASE 2의 표본 확장은 응답 감소를 해석하는 구성이 달라짐을 보여 준다. 이 결과는 H3가 제기한 비교 범위의 중요성과 부합한다. 다만 영향력 진단, 관측기간에 따른 표본 선택, 학습 분할은 서로 다른 문제이며 인과식별을 제공하지 않는다. 절차 응답의 변화가 실제 운영 변화인지도 완전히 분리하지 못했다.

\subsubsection{중심 연구 질문에 대한 관측자료 수준의 답변}
현재 자료로 주장할 수 있는 범위는 두 가지다. 첫째, CASE 1의 일부 지역 내 비교에서는 이주배경 인구비중과 평균 조사기간이 함께 증가하지만, 단면·차분·영향력 점검을 종합하면 이를 지역 전반의 일관된 기간 증가로 일반화하기 어렵다. 둘째, CASE 2에서는 고비중 집단의 상대적으로 높은 절차 충족률과 2025년 이후 격차 축소가 관측되며, 새 DML에서도 단기 축소 방향은 반복된다. 다만 분할에 따른 추론 변동과 표본별 응답 변화의 차이를 함께 보고해야 한다.

따라서 두 사례는 이주배경 인구가 집중된 지방정부에서 아동보호 행정 대응이 일관되게 더 느리다는 주장을 충분히 뒷받침하지 않는다. 이 판단은 기간 차이나 행정상 불이익이 없다는 결론도 아니다. 조사기간의 지역 내 연관성, 절차 충족률의 상대 수준, 단기 격차의 변화를 각각 확인한 결과다. 인구구성이 그 차이를 초래했다거나 2024년 개혁이 격차를 바꾸었다는 인과 주장, 개별 이주 아동이 더 늦게 보호받거나 등교했다는 주장은 현재의 집계자료로 식별할 수 없다.

\FloatBarrier
\begingroup\small\begin{thebibliography}{99}
\bibitem{case1} 제공 원고. \textit{이주배경 인구구성과 아동보호 행정의 조사기간}. Sweden\_ChildProtection\_Policy\_Draft\_KO.pdf, 26쪽. CASE 1 수치 및 기존 계산표.
\bibitem{case2} 제공 원고 (2026.11.11). \textit{집계된 개선의 이면: 스웨덴 보호시설 동반 아동의 학업보장 절차와 시군별 응답 전환}. Sweden\_VIND28\_Argument\_Revised\_20261111.pdf, 18쪽. CASE 2 원고 보고치.
\bibitem{reform} Socialstyrelsen (2024). \textit{Nya bestämmelser om skyddat boende}. Meddelandeblad 2/2024, 2024-3-8989.
\bibitem{kolada} RKA/Kolada. U33401 및 N02926 지표 메타데이터. \url{https://api.kolada.se/v3/kpi/U33401,N02926}.
\bibitem{scb} Statistiska centralbyrån. TAB4824 및 대도시권 분류. CASE 1 저장 응답과 제공 원고의 자료 정의.
\bibitem{vind} Socialstyrelsen (2026). \textit{Öppna jämförelser 2026 -- Våld i nära relationer}. Artikel 2026-6-10349. DATA·Analysstöd·Versionshistorik. 제공 CASE 2 원고의 출처. 이번에 공식 Excel DATA를 새로 확보했으나 당시 조사표 전체의 동등성을 재인증한 것은 아님.
\bibitem{sz} Sant'Anna, P. H. C. and Zhao, J. (2020). Doubly Robust Difference-in-Differences Estimators. \textit{Journal of Econometrics}, 219(1), 101--122. DOI: 10.1016/j.jeconom.2020.06.003.
\bibitem{dml} Chernozhukov, V. et al. (2018). Double/debiased Machine Learning for Treatment and Structural Parameters. \textit{The Econometrics Journal}, 21(1), C1--C68. DOI: 10.1111/ectj.12097.
\bibitem{sdid} Arkhangelsky, D. et al. (2021). Synthetic Difference-in-Differences. \textit{American Economic Review}, 111(12), 4088--4118. DOI: 10.1257/aer.20190159.
\bibitem{trygged} Trygged, S. (2017). Open Comparisons of Social Services in Sweden---Why, How, and for What? \textit{Cogent Social Sciences}, 3(1), 1404735. DOI: 10.1080/23311886.2017.1404735.
\bibitem{denvall} Denvall, V. et al. (2021). Lost in Comparison: A Program Theory Analysis of Performance Measurement in Social Work. \textit{Nordic Social Work Research}, 11(3), 249--263. DOI: 10.1080/2156857X.2019.1690559.

\bibitem{dmlscore} DoubleML. Score functions: observational panel difference-in-differences. \url{https://docs.doubleml.org/stable/guide/scores.html}. 2026년 11월 18일 확인. 본고는 공개 점수를 직접 구현하였다.
\bibitem{newraw} Socialstyrelsen. \textit{Öppna jämförelser 2026 -- Våld i nära relationer}, DATA 시트. \url{https://www.socialstyrelsen.se/contentassets/11a181f7bb3f4fd3b827c842b7df21c9/2026-6-10349.xlsx}. SCB TAB4823 및 TAB1121, 2015년. 2026년 11월 18일 공식 원자료 다운로드; 요청·해시는 동봉 재현자료 참조.
\end{thebibliography}\endgroup
\end{document}
